\section*{Chapter \ref{ch:b2:catalytic-cycles} --- Organometallic Catalysis: Elementary Steps and Cycles}

\begin{solution}{exo:b2:catalytic-cycles:1}
Before: Ir(I), $d^8$, $8 + 4 \times 2 = 16$ electrons. After: Ir(III), $d^6$, $6 + 6 \times 2 = 18$
electrons, octahedral.
\end{solution}

\begin{solution}{exo:b2:catalytic-cycles:2}
The first is a \omterm{def:b2:catalytic-cycles:oxidative-addition}{reductive elimination} (Pt(IV) to Pt(II), ethane formed from two methyls).
The second is a \omterm{def:b2:catalytic-cycles:ligand-exchange}{ligand exchange} (one phosphine replaced by ethene) followed by a
\omterm{def:b2:catalytic-cycles:insertion}{migratory insertion} of ethene into the Pd--Ph bond.
\end{solution}

\begin{solution}{exo:b2:catalytic-cycles:3}
TON $= 15.0/0.020 = 750$; TOF $= 750/3.0 = \qty{250}{h^{-1}}$.
\end{solution}

\begin{solution}{exo:b2:catalytic-cycles:4}
\begin{align*}
  &\mathrm{PdL_2 + ArBr \to ArPdBrL_2}\\
  &\mathrm{ArPdBrL_2 + Ar'B(OH)_3^- \to ArPdAr'L_2 + Br^- + B(OH)_3}\\
  &\mathrm{ArPdAr'L_2 \to Ar{-}Ar' + PdL_2}
\end{align*}
Sum: $\mathrm{ArBr + Ar'B(OH)_3^- \to Ar{-}Ar' + Br^- + B(OH)_3}$; every palladium species
cancels.
\end{solution}

\begin{solution}{exo:b2:catalytic-cycles:5}
Ethyl and isopropyl (they have $\beta$ hydrogens). Methyl has no $\beta$ carbon; neopentyl's $\beta$
carbon carries no hydrogen; benzyl's $\beta$ carbon is \omterm{def:b2:huckel:aromatic}{aromatic} and its hydrogen cannot reach
the metal.
\end{solution}

\begin{solution}{exo:b2:catalytic-cycles:6}
Methyl (\textit{E})-3-phenylpropenoate, \ce{Ph-CH=CH-COOCH3}, the phenyl added to the
terminal carbon, trans product.
\end{solution}

\begin{solution}{exo:b2:catalytic-cycles:7}
Butanal \ce{CH3CH2CH2CHO} (linear: the metal adds to the terminal carbon, the hydride to
the inner one) and 2-methylpropanal \ce{(CH3)2CHCHO} (branched: the metal on the inner
carbon).
\end{solution}

\begin{solution}{exo:b2:catalytic-cycles:8}
\ce{[RhCl(PPh3)3]} has 16 electrons: adding \ce{H2} (+2) would give 18 with seven
coordination positions, too crowded with three bulky phosphines. \ce{[RhCl(PPh3)2]} has 14
and an open site: \omterm{def:b2:catalytic-cycles:oxidative-addition}{oxidative addition} gives a 16-electron complex easily.
\end{solution}

\begin{solution}{exo:b2:catalytic-cycles:9}
The boronic acid is a poor nucleophile; the base adds hydroxide (or alkoxide) to boron,
giving a borate \ce{ArB(OH)3-}, whose aryl group is more electron-rich and transfers to
palladium.
\end{solution}

\begin{solution}{exo:b2:catalytic-cycles:10}
A 20-electron complex is not reasonable. The complex must first lose a ligand (14 e), then
add \ce{H2} (16 e), bind the alkene (18 e), insert it into an M--H bond (16 e) and eliminate the
alkane (14 e), before taking the ligand back or starting again.
\end{solution}

\begin{solution}{exo:b2:catalytic-cycles:11}
Initial rate $n_{\max}/\tau = \qty{0.25}{mmol/min}$; initial TOF $= 0.25/0.010 =
\qty{25}{min^{-1}} = \qty{1.5e3}{h^{-1}}$. Final TON $= 10.0/0.010 = 1.0 \times 10^3$.
\end{solution}

\begin{solution}{exo:b2:catalytic-cycles:12}
4-Bromoanisole with phenylboronic acid (or bromobenzene with 4-methoxyphenylboronic acid),
a palladium(0) phosphine catalyst and a base. Cycle: \omterm{def:b2:catalytic-cycles:oxidative-addition}{oxidative addition} of the
aryl bromide, \omterm{def:b2:catalytic-cycles:transmetalation}{transmetalation} from the borate, \omterm{def:b2:catalytic-cycles:oxidative-addition}{reductive elimination} of
4-methoxybiphenyl.
\end{solution}

\begin{solution}{pb:b2:catalytic-cycles:1}
\textbf{1.} Pd(II), $d^8$.
\textbf{2.} Pd(0), $d^{10}$, $10 + 2 \times 2 = 14$ electrons.
\textbf{3.} With 14 electrons it can accept two more pairs: it has \omterm{def:b2:catalytic-cycles:unsaturated}{vacant sites}.
\textbf{4.} Palladium(0) phosphine complexes are oxidised by air, and phosphines are
oxidised to phosphine oxides.
\textbf{5.} It reduces palladium(II) to palladium(0) (being oxidised itself) and binds the
palladium(0), keeping it in solution.
\textbf{6.} The \omterm{def:b2:catalytic-cycles:cycle}{precatalyst}.
\textbf{7.} \ce{Pd(PPh3)2 + CH3C6H4Br -> [Pd(C6H4CH3)Br(PPh3)2]}: Pd(II), $d^8$, $8 + 4 \times 2 =
16$ electrons.
\textbf{8.} It turns it into the borate \ce{PhB(OH)3-}, which transfers its phenyl group.
\textbf{9.} \ce{[Pd(Ar)Br(PPh3)2] + PhB(OH)3- -> [Pd(Ar)(Ph)(PPh3)2] + Br- + B(OH)3}.
\textbf{10.} \omterm{def:b2:catalytic-cycles:oxidative-addition}{Reductive elimination} joins two cis ligands; trans aryls must first isomerise.
\textbf{11.} \ce{[Pd(Ar)(Ph)(PPh3)2] -> Ar-Ph + Pd(PPh3)2}: palladium(0), 14 electrons, ready
for another turn.
\textbf{12.} \ce{CH3C6H4Br + PhB(OH)3- -> CH3C6H4-Ph + Br- + B(OH)3}.
\textbf{13.} 4-Bromotoluene (\qty{10.0}{mmol}, against \qty{12.0}{mmol} of boronic acid).
\textbf{14.} $1.51/168.24 = \qty{8.98e-3}{mol}$, \qty{8.98}{mmol}.
\textbf{15.} $8.98/10.0 = 89.8~\%$.
\textbf{16.} $0.050/10.0 = 0.50$~mol~\%.
\textbf{17.} $8.98/0.050 = 180$.
\textbf{18.} $180/4.0 = \qty{45}{h^{-1}}$.
\textbf{19.} Two phenyl groups transferred to the same palladium (homocoupling of the
boronic acid, favoured by traces of oxygen), then eliminated together.
\textbf{20.} The aryl groups have no $\beta$ hydrogen able to reach the metal.
\textbf{21.} The \omterm{def:b2:catalytic-cycles:oxidative-addition}{oxidative addition}: the \ce{C-Cl} bond is stronger than \ce{C-Br}.
\textbf{22.} It ends as palladium(0) particles or salts in the residues; it is expensive and
toxic, and is recovered.
\textbf{23.} Part of the boronic acid is lost to side reactions (homocoupling, loss of boron
from the ring); the excess keeps the bromide limiting.
\textbf{24.} The palladium turned over $\boldsymbol{\approx \num{1.8e2}}$ times.
\end{solution}
